\clearpage
\paragraph{\large Concluzii:}\mbox{}\\

\subparagraph{}{
Am formulat două probleme de programare liniară una de transport și una de producție, am încercat ca datele din aceste probleme să fie cât mai realiste și m-am asigurat că regiunea admisibilă nu este vidă.
}

\subparagraph{}{
Am utilizat metoda grafică pentru a rezolva o problema de programare liniară de transport și o problemă de producție. Problema de transport a fost puțin mai dificilă fiind că inițial avea șase variabile.
}


\subparagraph{}{
Uitândune mai bine la regiunile admisibile putem observa ca pentru problema de transport regiunea admisibilă este un poligon convex, iar in cazul în care poligonul nu ar fi convex problema nu ar putea fi rezolvata, pentru problema de transport regiunea admisibilă conturează un tronson.
Pentru ambele probleme obținem regiunea admisibilă mărginită.
Rezultatele acestor probleme sunt unice, soluții mai bune fiind imposibil de găsit.
}

\subparagraph{}{
O problemă care ar fi putut apărea în urma rezolvării acestor probleme este obținerea unor valori cu zecimale pentru variabile fiind că pare specific că la producție so produc unități întregi iar la transport nu este garantat că tranportul unei cantități înjumătățite ar reduce prețul de transport. 
}

\subparagraph{}{
	\LaTeX
}
